% !TEX root = main.tex
Dada la configuración $f^{3}$ estudiemos los términos j-j que podemos obtener. En primer lugar, $j \in \qty{5/2, 7/2}$ y por tanto, las combinaciones de $j$ para los 3 electrones son el conjunto: 
\begin{align*}
\qty(j_{1}, j_{2}, j_{3}) \in \left\{\qty(5/2, 5/2, 5/2),\;  \qty(5/2, 5/2, 7/2),\;  \qty(5/2, 7/2, 7/2),\;  \qty(7/2, 7/2, 7/2)\right\}
\end{align*}
La degeneración de cada combinación es: 
\begin{align*}
  g\qty(5/2, 5/2, 5/2) &= \binom{6}{3} = 20 &
g\qty(5/2, 5/2, 7/2) &= \binom{6}{2} \times  \binom{8}{1} = 120 \\
g\qty(5/2, 7/2, 7/2) &= \binom{6}{1} \times  \binom{8}{2} = 168 & 
g\qty(7/2, 7/2, 7/2) &= \binom{8}{3} = 56
\end{align*}
Y, la degeneración total es: 364. Resultado equivalente a $\binom{14}{3}$ como era de esperar. \par 
En este tipo de ejercicios, se suele preguntar por los términos j-j de la combinación de mayor degeneración. En este caso es la combinación $\qty(5/2, 7/2, 7/2)$. En esta combinación tenemos 2 electrones equivalentes con $j = 7/2$ y un electrón inequivalente con $j = 5/2$. Para tratarlo, primero acoplamos los electrones equivalentes, llamando $M_{J23}$ a la tercera componente de estos electrones acoplados. 
\begin{align*}
 \begin{array}{c | c | c}
   2M_{J23} & \text{ Particiones } & \text{ Términos }\\
   \hline 
   2M_{J23} = 12 & 12 = 7 + 5  & (7/2, 7/2)_{6} \\ 
   \hline 
    2M_{J23} = 10 & 10 = 7 + 3 & \text{ No hay nuevo término }\\ 
    \hline 
    2M_{J23} = 8 & \begin{aligned}
     8 &= 7 + 1\\
       &= 5 + 3
   \end{aligned} & (7/2, 7/2)_{4} \\
   \hline 
    2M_{J23} = 6 & \begin{aligned}
      6 &= 7 - 1\\
        &= 5 + 1\\
        \end{aligned} & \text{ No hay nuevo término }\\ 
        \hline
    2M_{J23} = 4 & \begin{aligned}
      4 &= 7 - 3\\
        &= 5 - 1\\
        &= 3 + 1
        \end{aligned}
     & (7/2, 7/2)_{2} \\
    \hline 
    2M_{J23} = 2 & \begin{aligned}
      2 &= 7 - 5\\ 
        &= 5 - 3\\
        &= 3 - 1\\ 
        \end{aligned} & \text{ No hay nuevo término }\\
        \hline 
    2M_{J23} = 0 & \begin{aligned}
      0 &= 7 - 7\\
        &= 5 - 5\\
        &= 3 - 3\\
        &= 1 - 1
    \end{aligned} & (7/2, 7/2)_{0}
\end{array}
    \end{align*}
    Donde recordemos la notación: $(j_{2},j_{3})_{J_{23}}$ . A continuación, acoplamos el electrón inequivalente con $j = 5/2$. Para ello simplemente descoponemos: 
    \begin{align*}
      J &= J_{23} \otimes j_{1} \implies \begin{cases}
        J = 6 \otimes 7/2 = 19/2 \oplus 17/2 \oplus 15/2 \oplus 13/2 \oplus 11/2 \oplus 9/2\\
      J = 4 \otimes 7/2 = 15/2 \oplus 13/2 \oplus 11/2 \oplus 9/2 \oplus 7/2 \oplus 5/2\\
      J = 2 \otimes 7/2 = 9/2 \oplus 7/2 \oplus 5/2 \oplus 3/2 \oplus 1/2 \\ 
      J = 0 \otimes 7/2 = 7/2
      \end{cases} 
    \end{align*}
    Por tanto, en los términos j-j de esta combinación son: 
    \begin{align*}
      \left\{ &(5/2, 7/2, 7/2)_{19/2}, \; (5/2, 7/2, 7/2)_{17/2}, \; 2 \times (5/2, 7/2, 7/2)_{15/2}, \; 2 \times (5/2, 7/2, 7/2)_{13/2},\right. \\ 
      &\; 2 \times (5/2, 7/2, 7/2)_{11/2}, \; 3 \times (5/2, 7/2, 7/2)_{9/2}, \; 3 \times (5/2, 7/2, 7/2)_{7/2}, \; 2 \times (5/2, 7/2, 7/2)_{5/2},\\ 
      &\left. \; (5/2, 7/2, 7/2)_{3/2}, \; (5/2, 7/2, 7/2)_{1/2} \right\}
    \end{align*}
    O en una notación más compacta: 
    \begin{align*}
      \boxed{
      (5/2, 7/2, 7/2)_\qty{19/2,\; 17/2,\;  2 \times 15/2,\;  2 \times 13/2,\;  2 \times 11/2, \; 3 \times 9/2,\;  3 \times 7/2,\;  2 \times 5/2,\;  3/2,\;  1/2}}
    \end{align*}
    \textbf{Nota:} En este caso, para calcular los términos j-j de los 2 $e^{-}$ es más sencillo estudiar la paridad de intercambio de los coeficientes de Clebsch-Gordan que acoplan $j_{2}$ y $j_{3}$ a $J_{23}$. Esta paridad viene dada por: 
    \begin{align*}
      P_{23} = (-1)^{J_{23} - j_{2} - j_{3}} = (-1)^{J_{23} - 1}
    \end{align*}
  Como $0 \leq J_{23} \leq 7$ entonces, simplemente podemos evaluar la anterior paridad para cada $J_{23}$ posible y descartar los términos pares que no cumplen la condición de antisimetría de los fermiones. Es directo darse cuenta que dichos términos son aquellos con $J_{23}$ par. Este método es más directo y nos permite comprobar nuestro resultado, pero es solo aplicable para el caso de 2 $e^{-}$ equivalentes. \par 
  La degeneración acumulada de este término es: 
\begin{align*}
  g_{T} &= 20 + 18 + 16 + 2 \times 14 + 2 \times 12 + 3 \times 10 + 3 \times 8 + 2 \times 6 + 4 = 168
\end{align*}    
Que es el resultado esperado.\par 
Por otro lado, por completitud estudiemos el caso de la combinación  hola $(7/2, 7/2, 7/2)$ donde tenemos 3 electrones equivalentes. Procedemos de la misma forma (véase el cálculo adjunto) y validamos el resultado con la degeneración acumulada: 
    \begin{align*}
    g_{T} = 16 + 12 + 10 + 8 + 6 + 4 = 56
    \end{align*}
 Que corresponde con la esperada. Así que finalmente los términos j-j de esta combinación: 
 \begin{align*}
   \boxed{
   (7/2, 7/2, 7/2)_\qty{15/2,\; 11/2,\;  9/2,\;  7/2,\;  5/2, \; 3/2,\;  3 \times 7/2,\; 5/2,\;  3/2}}
 \end{align*}

\begin{align*}
 \begin{array}{c | c | c}
   2M_{J} & \text{ Particiones } & \text{ Términos }\\
   \hline 
   2M_{J} = 15 & 15 = 7 + 5 + 3  & (7/2, 7/2, 7/2)_{15/2} \\ 
   \hline 
    2M_{J} = 13 & 13 = 7 + 5 + 1 & \text{ No hay nuevo término }\\ 
    \hline 
    2M_{J} = 11 & \begin{aligned}
     11 &= 7 + 5 -1\\
       &= 7 + 3 + 1
   \end{aligned} & (7/2, 7/2, 7/2)_{11/2} \\
   \hline 
    2M_{J} = 9 & \begin{aligned}
      9 &= 7 + 5 -3\\
        &= 7 + 3 -1\\
        &= 5 + 3 + 1\\
        \end{aligned} & (7/2, 7/2, 7/2)_{9/2}\\ 
        \hline
    2M_{J} = 7 & \begin{aligned}
      7 &= 7 + 3 - 3\\
        &= 7 + 5 - 5\\
        &= 7 + 1 - 1\\
        &= 5 + 3 - 1\\
        \end{aligned}
     & (7/2, 7/2, 7/2)_{7/2} \\
    \hline 
    2M_{J} = 5 & \begin{aligned}
      5 &= 5 + 3 - 3\\ 
        &= 7 + 3 - 5\\ 
        &= 7 + 1 - 3\\ 
        &= 7 + 5 - 7\\ 
        &= 5 + 1 - 1\\ 
        \end{aligned} & (7/2,7/2,7/2)_{5/2}\\
        \hline 
    2M_{J} = 3 & \begin{aligned}
      3 &= 5 + 1 - 3\\ 
        &= 5 + 3 - 5\\ 
        &= 7 + 1 - 5\\ 
        &= 7 + 3 - 7\\ 
        &= 7 - 1 - 3\\ 
        &= 3 + 1 - 1\\ 
        \end{aligned} & (7/2,7/2,7/2)_{3/2}\\
        \hline 
    2M_{J} = 1 & \begin{aligned}
      1 &= 3 + 1 - 3\\ 
        &= 5 - 1 - 3\\ 
        &= 5 + 1 - 5\\ 
        &= 5 + 3 - 7\\ 
        &= 7 - 1 - 5\\
        &= 7 + 1 - 7\\
        \end{aligned} & \text{ No hay nuevo término }\\
        \hline 
\end{array}
    \end{align*}
\begin{minipage}{0.37\textwidth} 
 \end{minipage} 
